% ============================================================
\section{The counting-collapse bound}
\label{sec:theory}

The mathematics here is classical (Banach fixed point, Lipschitz
composition, elementary softmax bounds); its value is precision, not
novelty: it fixes exactly which premise the experiments must test, and
Section~\ref{sec:contraction} reports that the premise fails on every
decoder measured. The development is machine-verified (end of section).

\subsection{Setting}

A decoder rendering a phrase repeated $k$ times passes through
\emph{repetition boundary} states $s_0, s_1, \dots$, the state after
each copy. Let $(S,d)$ be a nonempty complete metric space (in practice
$\mathbb{R}^m$, Euclidean) and $F \colon S \to S$ the boundary-to-boundary
map, $s_n = F^{\,n} s_0$. Two hypotheses connect this to a real decoder:
a \emph{single} $F$ acts at every boundary (Assumption~\ref{as:autonomous},
licensed approximately by Lemma~\ref{lem:dilution}), and $F$ contracts
(the antecedent Section~\ref{sec:contraction} measures and rejects).

\subsection{Attention dilution}

\begin{lemma}[Attention dilution]
\label{lem:dilution}
Let $z \in \mathbb{R}^k$ be the attention logits of the $k$ text positions
occupied by the $k$ copies of the repeated phrase, and suppose their
spread is bounded: $z_a - z_b \le \delta$ for all $a, b$. Let
$p_i = e^{z_i} / \sum_{j=1}^{k} e^{z_j}$ be the attention weights. Then
\begin{enumerate}
\item[(i)] $\displaystyle \frac{e^{-\delta}}{k} \;\le\; p_i \;\le\;
      \frac{e^{\delta}}{k}$ \quad for every $i$;
\item[(ii)] $\displaystyle |p_i - p_j| \;\le\;
      \frac{e^{\delta} - e^{-\delta}}{k}$ \quad for every $i, j$;
\item[(iii)] $\displaystyle H(p) \;=\; -\sum_{i=1}^{k} p_i \log p_i
      \;\ge\; \log k - \delta$.
\end{enumerate}
\end{lemma}

\begin{proof}
(i) For every $j$, the spread bound gives
$e^{z_j} \ge e^{z_i - \delta}$, so
$\sum_j e^{z_j} \ge k\, e^{z_i - \delta}$ and
$p_i \le e^{z_i} / (k\, e^{z_i - \delta}) = e^{\delta}/k$. Symmetrically,
$e^{z_j} \le e^{z_i + \delta}$ gives
$\sum_j e^{z_j} \le k\, e^{z_i + \delta}$ and
$p_i \ge e^{-\delta}/k$.
(ii) is immediate from (i).
(iii) From (i), $\log p_i \le \delta - \log k$, so
$-\log p_i \ge \log k - \delta$; multiplying by $p_i \ge 0$ and summing
with $\sum_i p_i = 1$ gives the claim.
\end{proof}

The profile is within $(e^{\delta}-e^{-\delta})/k$ of uniform and its
entropy within $\delta$ of maximal: attention weight carries essentially
no information about \emph{which} occurrence is being rendered, licensing,
up to $O(\delta + 1/k)$:

\begin{assumption}[Autonomous boundary map]
\label{as:autonomous}
The same map $F$ carries the state from each repetition boundary to the
next.
\end{assumption}

\begin{remark}[What a formal bridge would need]
\label{rem:bridge}
The step from Lemma~\ref{lem:dilution} to
Assumption~\ref{as:autonomous} is quantitative but informal, and is
deliberately not part of the verified development. Making it formal would
require: a context-indexed family of maps $F_c$ with a metric on
contexts; a bound on the attended \emph{value} vectors, not only on the
attention weights; a perturbed-contraction (shadowing) lemma; and a
re-derivation of Theorem~\ref{thm:collapse} carrying an
$O(\kappa\rho/(1-q))$ correction, where $\rho$ bounds the per-boundary
perturbation of the conditioning and $\kappa$ the decoder's sensitivity
to it. The bridge remains, honestly, a remark.
\end{remark}

\subsection{The bound}

\begin{theorem}[Counting collapse]
\label{thm:collapse}
Suppose $F$ is a contraction with factor $q \in (0,1)$, that is,
$d(Fx, Fy) \le q\, d(x,y)$ for all $x, y \in S$. Fix $s_0 \in S$, let
$s^\ast$ be the unique fixed point of $F$, and set
\[
  C \;=\; \frac{d(s_0, F s_0)}{1 - q}.
\]
Then, for all $m, n \in \mathbb{N}$:
\begin{enumerate}
\item[(i)] \textup{(geometric convergence)}\quad
      $d(F^{\,n} s_0,\, s^\ast) \;\le\; C q^{\,n}$;
\item[(ii)] \textup{(eventual indistinguishability)}\quad
      $d(F^{\,m} s_0,\, F^{\,n} s_0) \;\le\; C\,(q^{\,m} + q^{\,n})$, so
      for every $\varepsilon > 0$ there is an $N$ with
      $d(F^{\,m} s_0, F^{\,n} s_0) < \varepsilon$ whenever $m, n \ge N$;
\item[(iii)] \textup{(no Lipschitz readout can count)}\quad for every
      $L$-Lipschitz $g \colon S \to \mathbb{R}$,
      \[
        \bigl|\, g(F^{\,m} s_0) - g(F^{\,n} s_0) \,\bigr|
        \;\le\; L\, C\, (q^{\,m} + q^{\,n});
      \]
\item[(iv)] \textup{(finite counting horizon)}\quad if $m \le n$ and a
      readout still separates the two counts with decision margin
      $\mu > 0$, that is
      $\mu \le |g(F^{\,m} s_0) - g(F^{\,n} s_0)|$, then
      $\mu \le 2 L C q^{\,m}$ and hence
      \[
        m \;\le\; \frac{\log\bigl(\mu / (2 L C)\bigr)}{\log q}
          \;=\; \frac{\log\bigl(2 L C / \mu\bigr)}{\log (1/q)}
          \;=:\; N^\ast .
      \]
\end{enumerate}
\end{theorem}

\begin{proof}
(i) By the triangle inequality and the contraction property,
$d(s_0, s^\ast) \le \sum_{j \ge 0} d(F^{\,j} s_0, F^{\,j+1} s_0)
\le \sum_{j \ge 0} q^{\,j}\, d(s_0, F s_0) = C$. Applying $F$ $n$ times,
$d(F^{\,n} s_0, s^\ast) = d(F^{\,n} s_0, F^{\,n} s^\ast)
\le q^{\,n} d(s_0, s^\ast) \le C q^{\,n}$.

(ii) Triangle inequality through the fixed point:
$d(F^{\,m} s_0, F^{\,n} s_0) \le d(F^{\,m} s_0, s^\ast) +
d(F^{\,n} s_0, s^\ast) \le C(q^{\,m} + q^{\,n})$. Given
$\varepsilon > 0$, choose $N$ with
$q^{\,N} < \varepsilon / (2C + 1)$; then for $m, n \ge N$ the bound is
below $\varepsilon$.

(iii) $|g(x) - g(y)| \le L\, d(x,y)$ composed with (ii).

(iv) Since $q < 1$ and $m \le n$, $q^{\,n} \le q^{\,m}$, so (iii) gives
$\mu \le 2 L C q^{\,m}$. The prefactor $2LC$ is strictly positive
(otherwise the bound would force $\mu \le 0$), so
$\mu / (2LC) \le q^{\,m}$; taking logarithms, and dividing by
$\log q < 0$, reverses the inequality and yields the horizon.
\end{proof}

The readout $g$ is any Lipschitz function of the state: the stop-token
logit, a duration predictor, any linear probe. Beyond $N^\ast$ the decoder
provably cannot tell how many repetitions it has produced. The implication
is genuine; its antecedent is not (measured $q$ between 21.0 and 347.7,
never below 1; Section~\ref{sec:contraction}).

\subsection{Machine verification}

Lemma~\ref{lem:dilution}(i)--(iii) and
Theorem~\ref{thm:collapse}(i)--(iv), with the auxiliary inequalities
their proofs use, are verified in Lean~4 over mathlib: no unproven
placeholders, and an axiom audit shows every exported theorem depends only
on the three standard axioms (propositional extensionality, choice,
quotient soundness). The verified statements match the theorems printed
here one-to-one; the only deliberately unverified step is the bridge of
Remark~\ref{rem:bridge}, informal there exactly as here. The development
ships with the release (Section~\ref{sec:release}).
